\documentclass[12pt, a4paper]{article}

% --- Pacotes Fundamentais ---
\usepackage[utf8]{inputenc}
\usepackage[T1]{fontenc}
\usepackage[brazil]{babel}
\usepackage{geometry}
\usepackage{titlesec}
\usepackage{fancyhdr}
\usepackage{graphicx}
\usepackage{xcolor}
\usepackage{booktabs}
\usepackage{longtable}
\usepackage{tabularx}
\usepackage{colortbl}
\usepackage{float}
\usepackage{parskip}
\usepackage{lmodern}
\usepackage{lastpage}
\usepackage{tikz}
\usepackage{pgfplots}
\usepackage{amssymb}
\pgfplotsset{compat=1.18}

% --- Configurações de Layout ---
\geometry{left=2.5cm, right=2.5cm, top=4cm, bottom=2.5cm, headsep=1cm, headheight=60pt}

% --- Definição de Cores ---
\definecolor{primaryColor}{RGB}{0, 51, 102}   % Azul Marinho
\definecolor{secondaryText}{RGB}{80, 80, 80}  % Cinza Escuro
\definecolor{accentColor}{RGB}{178, 34, 34}   % Vermelho Tijolo
\definecolor{lightGray}{RGB}{245, 245, 245}
\definecolor{tableHeader}{RGB}{230, 230, 250}
\definecolor{successGreen}{RGB}{34, 139, 34}  % Verde para ROAS positivo
\definecolor{warningOrange}{RGB}{255, 140, 0} % Laranja para validação

% --- Configuração de Cabeçalho e Rodapé ---
\pagestyle{fancy}
\fancyhf{}

% Cabeçalho
\fancyhead[C]{%
    \begin{tabularx}{\textwidth}{@{} X r @{}}
        \textcolor{primaryColor}{\textbf{\footnotesize FRAMEWORK DE ESCALA: TRÁFEGO PAGO LOW TICKET 2026}} & \scriptsize \textit{Metodologia MVT} \\
        \textcolor{secondaryText}{\scriptsize Modelagem, Vantagem e Tráfego} & {\scriptsize \textbf{Data Base: 07/01/2026}}
    \end{tabularx}%
}

% Rodapé
\fancyfoot[R]{\normalsize Página \textbf{\thepage} de \textbf{\pageref{LastPage}}}

% --- Linha separadora do CABEÇALHO ---
\renewcommand{\headrulewidth}{1.5pt}
\renewcommand{\headrule}{%
    \vspace{0.1cm}
    \hbox to\headwidth{\color{primaryColor}\leaders\hrule height \headrulewidth\hfill}%
}

% --- Linha separadora do RODAPÉ ---
\renewcommand{\footrulewidth}{1.5pt}
\renewcommand{\footrule}{%
    \color{primaryColor}%
    \hbox to\headwidth{\leaders\hrule height \footrulewidth\hfill}%
    \vspace{0.2cm}
}

% --- Formatação de Títulos ---
\titleformat{\section}
{\color{primaryColor}\normalfont\Large\bfseries}
{\thesection}{1em}{}[{\titlerule[0.8pt]}]

\titleformat{\subsection}
{\color{primaryColor}\normalfont\large\bfseries}
{\thesubsection}{1em}{}

% --- Macro para Moeda ---
\newcommand{\reais}[1]{R\$ #1}

% --- INÍCIO DO DOCUMENTO ---
\begin{document}

% --- CAPA ---
\begin{titlepage}
    \centering
    \vspace*{3cm}
    {\Huge \textbf{\textcolor{primaryColor}{FRAMEWORK DE ESCALA}}} \\
    \vspace{0.5cm}
    {\Large \textbf{TRÁFEGO PAGO LOW TICKET 2026}} \\
    \vspace{0.3cm}
    {\large Metodologia MVT e Ciclo de Vida de Ofertas no Facebook Ads} \\
    \vspace{2cm}
    \rule{10cm}{1.5pt} \\
    \vspace{0.5cm}
    \colorbox{lightGray}{%
        \parbox{0.9\textwidth}{%
            \centering
            \vspace{0.3cm}
            {\Large \textbf{GUIA ESTRATÉGICO DE CRESCIMENTO}} \\
            \vspace{0.3cm}
            {\large Como escalar de \reais{50,00} para \reais{500,00} diários} \\
            \vspace{0.3cm}
        }%
    }
    \vspace{2cm}

    \begin{flushleft}
    \textbf{PILARES DA METODOLOGIA:} \\
    \vspace{0.2cm}
    \textcolor{primaryColor}{\textbf{M}} - Modelagem de Ofertas Validadas \\
    \textcolor{primaryColor}{\textbf{V}} - Vantagem Competitiva (Marketing de Nicho) \\
    \textcolor{primaryColor}{\textbf{T}} - Tráfego Pago Estruturado e Escalável
    \end{flushleft}

    \vfill
    {\small Framework prático para produtos digitais de baixo ticket, contemplando as novas regras fiscais do Facebook (Imposto de 12,5\% sobre investimento em tráfego) e estruturas validadas em campanhas reais rodando em dezembro/2025 e janeiro/2026.} \\
    \vspace{1cm}
    \textbf{Data Base:} 07 de Janeiro de 2026
\end{titlepage}

% --- CONTEÚDO ---
\section{O Cenário do Tráfego Pago em 2026}

O ano de 2026 trouxe mudanças estruturais no Facebook Ads que impactam diretamente a rentabilidade de campanhas low ticket. A principal alteração é a \textbf{tributação de 12,5\% sobre todo o investimento em tráfego}, independentemente do resultado da campanha.

\subsection{Impacto da Nova Legislação}

\begin{table}[H]
    \centering
    \renewcommand{\arraystretch}{1.4}
    \begin{tabularx}{\textwidth}{X c r}
        \toprule
        \rowcolor{tableHeader} \textbf{Cenário de Investimento} & \textbf{Verba} & \textbf{Custo Tributário} \\
        \midrule
        Investimento Diário Básico & \reais{100,00} & \reais{12,50} \\
        Investimento Escala Média & \reais{500,00} & \reais{62,50} \\
        Investimento Escala Alta & \reais{1.000,00} & \reais{125,00} \\
        \midrule
        \rowcolor{accentColor!10} \textbf{Investimento Mensal (R\$ 1.000/dia)} & \textbf{\reais{30.000,00}} & \textbf{\reais{3.750,00}} \\
        \bottomrule
    \end{tabularx}
    \caption*{\footnotesize \textbf{Atenção:} Este custo tributário incide mesmo em campanhas com ROAS 0x0 ou negativo, tornando essencial a validação de estruturas antes da escala.}
\end{table}

\subsection{Consequências Estratégicas}

Esta nova realidade elimina do mercado dois perfis de operadores:

\begin{enumerate}
    \item \textbf{Mineradores de Ofertas:} Aqueles que vasculham a Biblioteca de Anúncios copiando criativos e estruturas alheias sem compreender a metodologia subjacente enfrentarão prejuízos sistemáticos.

    \item \textbf{Testadores Compulsivos:} Operadores que sobem 10, 15, 20 criativos simultaneamente sem critérios de validação dilapidarão capital em testes desestruturados.
\end{enumerate}

\textcolor{accentColor}{\textbf{A solução:}} Dominar frameworks de validação progressiva, compreendendo as \textbf{6 Fases da Jornada Low Ticket} e aplicando estruturas de campanha com critérios objetivos de decisão.

\newpage

\section{As 6 Fases da Jornada Low Ticket}

O ciclo de vida de uma oferta low ticket no Facebook Ads é composto por seis fases distintas, cada uma com objetivos, métricas e tomadas de decisão específicas.

\begin{table}[H]
    \centering
    \small
    \renewcommand{\arraystretch}{1.5}
    \begin{tabularx}{\textwidth}{c X X}
        \toprule
        \rowcolor{tableHeader} \textbf{Fase} & \textbf{Objetivo Principal} & \textbf{Duração Típica} \\
        \midrule
        \textbf{1} & \textbf{Teste Criativo:} Validar criativos com ROAS $\geq$ 1.8 & 2-3 dias \\
        \rowcolor{lightGray}
        \textbf{2} & \textbf{Escala:} Aumentar orçamento em estruturas validadas & 3-7 dias \\
        \textbf{3} & \textbf{Teto de Oferta:} Identificar limite de saturação da demanda & Variável \\
        \rowcolor{lightGray}
        \textbf{4} & \textbf{Análise de Padrões:} Mapear estruturas de melhor performance & Contínuo \\
        \textbf{5} & \textbf{Manutenção:} Garantir ROAS consistente sem intervenções diárias & 7-15 dias \\
        \rowcolor{lightGray}
        \textbf{6} & \textbf{Saturação:} Reconhecer queda irreversível e encerrar oferta & 3-5 dias \\
        \bottomrule
    \end{tabularx}
\end{table}

\subsection{Princípio Fundamental}

\begin{quote}
\textcolor{primaryColor}{\textbf{``Não se escala número de vendas. Escala-se ROAS e lucro.''}}
\end{quote}

O erro crítico de gestores inexperientes é aumentar orçamento baseado em volume de conversões (``Este criativo vendeu 10 vezes, vou duplicar a verba''). A métrica correta de decisão é o \textbf{Retorno sobre Investimento Publicitário (ROAS)}, que indica a sustentabilidade financeira da operação.

\textcolor{accentColor}{\textbf{Regra de Ouro:}} Somente criativos com ROAS mínimo de 1.8 na fase de teste justificam progressão para escala. Criativos com ROAS entre 1.3 e 1.5 permanecem em observação. Abaixo de 1.3, desativação imediata.

\newpage

\section{Protocolo de Teste de Criativo (Fase de Validação)}

A Fase 1 é a fundação de toda operação escalável. Seu objetivo não é gerar lucro imediato, mas sim \textbf{identificar ativos criativos com potencial de escala}.

\subsection{Estrutura ABO Validada (Ad Set Budget Optimization)}

\begin{table}[H]
    \centering
    \renewcommand{\arraystretch}{1.4}
    \begin{tabularx}{\textwidth}{X X}
        \toprule
        \rowcolor{tableHeader} \textbf{Elemento da Estrutura} & \textbf{Configuração} \\
        \midrule
        Número de Campanhas & 1 campanha por criativo testado \\
        \rowcolor{lightGray}
        Conjuntos por Campanha & 3 conjuntos idênticos (segmentação, idade, gênero) \\
        Anúncios por Conjunto & 1 criativo (mesmo em todos os conjuntos) \\
        \rowcolor{lightGray}
        Orçamento por Conjunto & 45\%-50\% do ticket médio do produto \\
        Público-Alvo & Advantage+ Audience (sugestão de gênero e idade) \\
        \rowcolor{lightGray}
        Tempo Mínimo de Teste & 2 dias (ideal: 3 dias para confirmação de constância) \\
        \bottomrule
    \end{tabularx}
\end{table}

\subsection{Lógica da Estrutura 1-3-1}

A arquitetura de \textbf{1 Campanha $\rightarrow$ 3 Conjuntos $\rightarrow$ 1 Anúncio} não testa públicos diferentes, mas sim \textbf{fatias distintas do mesmo público} processadas pelo algoritmo do Facebook. Isso permite:

\begin{itemize}
    \item Confirmar que o criativo tem apelo consistente (não é resultado de sorte em uma única fatia)
    \item Identificar variações de CPM (Custo por Mil Impressões) entre diferentes segmentos de entrega
    \item Validar estabilidade de métricas secundárias (CPV, taxa de checkout)
\end{itemize}

\subsection{Exemplo de Orçamento}

Para um produto com ticket médio de \reais{37,00}:

\begin{itemize}
    \item \textbf{Orçamento por Conjunto:} \reais{17,00} a \reais{20,00} (45\%-50\% do ticket)
    \item \textbf{Investimento Total por Criativo:} \reais{51,00} a \reais{60,00} (3 conjuntos)
    \item \textbf{Teste de 3 criativos simultâneos:} \reais{153,00} a \reais{180,00}
\end{itemize}

\textcolor{warningOrange}{\textbf{Estrutura Simplificada:}} Para orçamentos mais restritos, é aceitável usar \textbf{2 conjuntos por criativo} com orçamento mínimo de 35\% do ticket.

\newpage

\subsection{KPIs de Sucesso no Teste de Criativo}

Um criativo é considerado \textbf{validado} quando atende simultaneamente aos seguintes critérios:

\begin{table}[H]
    \centering
    \renewcommand{\arraystretch}{1.5}
    \begin{tabularx}{\textwidth}{X c c}
        \toprule
        \rowcolor{tableHeader} \textbf{Métrica} & \textbf{Range Ideal} & \textbf{Status} \\
        \midrule
        \textbf{ROAS (Retorno sobre Gasto)} & $\geq$ 1.8 & \textcolor{successGreen}{VALIDADO} \\
        \rowcolor{lightGray}
        ROAS Zona de Observação & 1.5 - 1.8 & \textcolor{warningOrange}{MANTER ATIVO} \\
        ROAS Abaixo do Mínimo & $<$ 1.3 & \textcolor{accentColor}{DESATIVAR} \\
        \midrule
        \textbf{CPV (Custo por Visualização)} & 5\% - 6\% & Métrica Secundária \\
        \rowcolor{lightGray}
        \textbf{Taxa de Checkout} & 25\% - 30\% & Métrica Secundária \\
        \textbf{CPM (Custo por Mil Impressões)} & \reais{15,00} - \reais{25,00} & Métrica Secundária \\
        \bottomrule
    \end{tabularx}
\end{table}

\subsection{Importância das Métricas Secundárias}

Enquanto o ROAS é o indicador primário de viabilidade econômica, as métricas secundárias revelam a \textbf{saúde do funil de conversão}:

\begin{itemize}
    \item \textbf{CPV baixo (5\%-6\%):} Indica que o criativo está gerando tráfego qualificado para a página de vendas a um custo eficiente.
    \item \textbf{Taxa de Checkout elevada (25\%-30\%):} Confirma que a página de vendas possui copy persuasivo e alinhamento com a promessa do anúncio.
    \item \textbf{CPM dentro da faixa:} Sinaliza que o Facebook não está cobrando prêmio por segmentação excessiva (sinal de público adequado).
\end{itemize}

\textcolor{accentColor}{\textbf{Atenção:}} Criativos com ROAS alto mas métricas secundárias ruins tendem a ter \textbf{baixa durabilidade na escala}. Priorize criativos que atendam todos os critérios.

\newpage

\section{Metodologia de Escala: De \reais{50,00} para \reais{500,00}+}

Após validar criativos na Fase 1, inicia-se a Fase 2: \textbf{Escala Progressiva}. O objetivo é aumentar o investimento diário mantendo ou melhorando o ROAS.

\subsection{Princípio da Escala Validada}

\begin{quote}
\textcolor{primaryColor}{\textbf{``Escala não vem do nada. Ela vem de algo validado.''}}
\end{quote}

Jamais escale campanhas com ROAS abaixo de 1.8. O aumento de orçamento em estruturas marginais (ROAS 1.3-1.5) invariavelmente resulta em ROAS ainda menor, transformando operações neutras em prejuízo.

\subsection{Duas Rotas de Escala}

Existem duas arquiteturas principais para escalar ofertas low ticket:

\begin{table}[H]
    \centering
    \small
    \renewcommand{\arraystretch}{1.4}
    \begin{tabularx}{\textwidth}{X X X}
        \toprule
        \rowcolor{tableHeader} \textbf{Modelo} & \textbf{Estrutura} & \textbf{Orçamento Recomendado} \\
        \midrule
        \textbf{Escala em ABO} & Duplicação de conjuntos validados com aumento de verba & 3x a 6x o ticket médio \\
        \rowcolor{lightGray}
        \textbf{Escala em CBO} & Campaign Budget Optimization (1-3-1) & 5x a 6x o ticket médio \\
        \bottomrule
    \end{tabularx}
\end{table}

\subsection{Escala em ABO: Modelo 01 (Duplicação)}

Quando um criativo performa bem no teste (ROAS $\geq$ 1.8 por 2-3 dias consecutivos):

\begin{enumerate}
    \item Duplique o conjunto vencedor 1 a 2 vezes \textbf{dentro da mesma campanha}
    \item Aumente o orçamento para 3x o ticket médio (ex: produto de \reais{37,00} $\rightarrow$ orçamento de \reais{110,00})
    \item Mantenha a estrutura ativa por 48-72h antes de avaliar novos movimentos
\end{enumerate}

\textcolor{warningOrange}{\textbf{Preferência Atual:}} Menos conjuntos com mais verba (concentração) ao invés de pulverização em múltiplas campanhas.

\subsection{Escala em ABO: Modelo 02 (Multi-Criativo)}

Agrupe 3 a 4 criativos validados em uma única campanha:

\begin{enumerate}
    \item Crie 1 campanha com 1 conjunto de anúncios
    \item Insira 3 criativos que performaram bem no teste (ROAS individual $\geq$ 1.8)
    \item Orçamento: 3x a 4x o ticket médio
    \item O Facebook distribuirá orçamento entre os criativos conforme performance
\end{enumerate}

\newpage

\subsection{Escala em CBO: Modelo 1-3-1}

Para criativos com \textbf{constância validada} (ROAS estável por 5-7 dias):

\begin{table}[H]
    \centering
    \renewcommand{\arraystretch}{1.4}
    \begin{tabularx}{\textwidth}{X X}
        \toprule
        \rowcolor{tableHeader} \textbf{Elemento} & \textbf{Configuração} \\
        \midrule
        Número de Campanhas & 1 (CBO ativado) \\
        \rowcolor{lightGray}
        Conjuntos de Anúncios & 3 conjuntos idênticos \\
        Anúncio por Conjunto & 1 criativo (mesmo em todos) \\
        \rowcolor{lightGray}
        Orçamento Total da Campanha & 5x a 6x o ticket médio \\
        Estratégia de Lance & Menor Custo (padrão) \\
        \bottomrule
    \end{tabularx}
\end{table}

\textbf{Exemplo Prático:}
\begin{itemize}
    \item Produto: \reais{37,00}
    \item Criativo 31 validado com ROAS 2.1 por 7 dias em ABO
    \item Orçamento CBO: \reais{200,00} (5.4x o ticket)
    \item Resultado esperado: ROAS entre 1.6 e 2.0 (queda aceitável vs. fase de teste)
\end{itemize}

\subsection{Gestão de Expectativa na Escala}

\textcolor{accentColor}{\textbf{Fenômeno Normal:}} Nos primeiros 1-3 dias da escala, é comum o ROAS ficar neutro (1.0x) ou até negativo (0.8x-0.9x). Isso ocorre porque:

\begin{itemize}
    \item O algoritmo está re-otimizando para o novo patamar de orçamento
    \item A entrega está sendo testada em novas fatias de público
    \item O CPM pode aumentar temporariamente
\end{itemize}

\textbf{Regra de Paciência:} Mantenha estruturas de escala ativas por no mínimo 3 dias completos antes de desativar, \textbf{desde que o ROAS acumulado não caia abaixo de 0.7x}. Muitas campanhas que iniciam ruins tornam-se lucrativas no 4º ou 5º dia.

\newpage

\section{Metodologia MVT: Modelagem, Vantagem e Tráfego}

A sustentabilidade de longo prazo no tráfego pago low ticket não depende apenas de habilidades técnicas no Facebook Ads, mas da \textbf{concepção estratégica da oferta}.

\subsection{O Problema da Cópia}

Operadores que mineram ofertas bem-sucedidas na Biblioteca de Anúncios enfrentam três desvantagens estruturais:

\begin{enumerate}
    \item \textbf{Saturação Prévia:} A oferta já foi exposta ao mercado por semanas/meses, esgotando demanda latente
    \item \textbf{Guerra de Lance:} Múltiplos anunciantes competindo pela mesma audiência elevam o CPM artificialmente
    \item \textbf{Custo de Tributação:} Com o imposto de 12,5\%, margens já comprimidas pela concorrência tornam-se insustentáveis
\end{enumerate}

\subsection{A Alternativa MVT}

\begin{table}[H]
    \centering
    \renewcommand{\arraystretch}{1.5}
    \begin{tabularx}{\textwidth}{c X}
        \toprule
        \rowcolor{tableHeader} \textbf{Pilar} & \textbf{Descrição} \\
        \midrule
        \textbf{M - Modelagem} & Estudar estruturas vencedoras do mercado não para copiar, mas para entender princípios de conversão (gatilhos mentais, arquitetura de oferta, prova social). \\
        \rowcolor{lightGray}
        \textbf{V - Vantagem} & Criar concepções únicas direcionadas a nichos específicos onde não há concorrência direta. Posicionamento de ``Oceano Azul''. \\
        \textbf{T - Tráfego} & Aplicar frameworks de escala validados (este documento) para maximizar ROI e sustentabilidade. \\
        \bottomrule
    \end{tabularx}
\end{table}

\subsection{Ciclo de Vida da Vantagem Competitiva}

\begin{center}
\begin{tikzpicture}
    \begin{axis}[
        width=14cm,
        height=8cm,
        xlabel={Tempo (semanas)},
        ylabel={Lucro Semanal (R\$)},
        xmin=0, xmax=12,
        ymin=0, ymax=5000,
        xtick={0,2,4,6,8,10,12},
        ytick={0,1000,2000,3000,4000,5000},
        legend pos=north west,
        ymajorgrids=true,
        grid style=dashed,
    ]

    \addplot[
        color=blue,
        mark=square,
        thick
        ]
        coordinates {
        (1,800)(2,1500)(3,2800)(4,4200)(5,4500)(6,4300)(7,3800)(8,3000)(9,2200)(10,1500)(11,900)(12,500)
        };
        \addlegendentry{Oferta com Vantagem (MVT)}

    \addplot[
        color=red,
        mark=triangle,
        thick,
        dashed
        ]
        coordinates {
        (1,200)(2,300)(3,400)(4,500)(5,450)(6,350)(7,250)(8,150)(9,100)(10,50)(11,0)(12,-100)
        };
        \addlegendentry{Oferta Copiada (Saturada)}

    \end{axis}
\end{tikzpicture}
\end{center}

\textbf{Interpretação:} Ofertas criadas com vantagem competitiva geram lucro expressivo nas primeiras 4-8 semanas, antes que a concorrência replique o posicionamento. Ofertas copiadas iniciam com lucro marginal e rapidamente tornam-se deficitárias.

\newpage

\section{Fluxo de Decisão Pós-Teste de Criativo}

Após completar 2-3 dias de teste com a estrutura ABO 1-3-1, você deve analisar os resultados e tomar uma das seguintes ações:

\begin{table}[H]
    \centering
    \small
    \renewcommand{\arraystretch}{1.4}
    \begin{tabularx}{\textwidth}{c X X}
        \toprule
        \rowcolor{tableHeader} \textbf{ROAS} & \textbf{Classificação} & \textbf{Ação Recomendada} \\
        \midrule
        \textcolor{successGreen}{$\geq$ 1.8} & \textbf{VALIDADO} & Levar para escala (Modelo ABO ou CBO) \\
        \rowcolor{lightGray}
        \textcolor{warningOrange}{1.5 - 1.8} & \textbf{EM VALIDAÇÃO} & Manter ativo por mais 1-2 dias \\
        \textcolor{accentColor}{1.3 - 1.5} & \textbf{MARGINAL} & Desativar se não melhorar em 24h \\
        \rowcolor{lightGray}
        \textcolor{accentColor}{$<$ 1.3} & \textbf{REPROVADO} & Pausar imediatamente \\
        \bottomrule
    \end{tabularx}
\end{table}

\subsection{Cenário 1: Nenhum Criativo Validado}

Se após testar 3-5 criativos todos apresentarem ROAS $<$ 1.3, realize os seguintes ajustes:

\begin{enumerate}
    \item \textbf{Troca de Estrutura:} Teste CBO ao invés de ABO (ou vice-versa)
    \item \textbf{Novos Criativos:} Produza variações com diferentes ângulos de copy/creative
    \item \textbf{Ajuste na Página:} Revise headline, prova social, garantia, call-to-action
    \item \textbf{Novo Checkout:} Teste processadores de pagamento alternativos (taxas menores)
    \item \textbf{Nova Oferta:} Em último caso, a concepção pode não ter fit com o mercado
\end{enumerate}

\subsection{Cenário 2: 1-2 Criativos Validados}

Este é o cenário ideal. Concentre recursos nos criativos vencedores:

\begin{enumerate}
    \item Pause todos os criativos com ROAS $<$ 1.5
    \item Mantenha criativos com ROAS 1.5-1.8 em observação (sem aumentar orçamento)
    \item Escale imediatamente os criativos com ROAS $\geq$ 1.8 usando Modelo ABO 01 ou CBO
\end{enumerate}

\subsection{Cenário 3: Múltiplos Criativos Validados (3+)}

Quando você tem 3 ou mais criativos com ROAS $\geq$ 1.8, você pode:

\begin{itemize}
    \item Criar uma campanha ABO Modelo 02 (top 3-4 criativos agrupados)
    \item Criar múltiplas campanhas CBO (1 criativo por campanha)
    \item Testar ambas as estruturas simultaneamente (cardápio de opções)
\end{itemize}

\textcolor{primaryColor}{\textbf{Filosofia do Cardápio:}} Não existe ``estrutura perfeita''. Ofertas diferentes respondem melhor a estruturas diferentes. Teste ABO e CBO em paralelo para identificar qual funciona melhor \textbf{para aquela oferta específica}.

\newpage

\section{Estudo de Caso Real: Dezembro/2025 - Janeiro/2026}

Esta seção documenta uma operação real executada com produto low ticket de \reais{37,00} (ticket médio), demonstrando a aplicação prática do framework.

\subsection{Linha do Tempo da Operação}

\begin{table}[H]
    \centering
    \footnotesize
    \renewcommand{\arraystretch}{1.3}
    \begin{tabularx}{\textwidth}{c X r r}
        \toprule
        \rowcolor{tableHeader} \textbf{Data} & \textbf{Fase/Ação} & \textbf{Investimento} & \textbf{ROAS} \\
        \midrule
        10-11/12 & Teste Criativo (6 conjuntos, 3 criativos) & \reais{250,00} & 1.48 \\
        \rowcolor{lightGray}
        12-13/12 & Escala Modelo ABO (Criativos 31, 6, 14) & \reais{520,00} & 1.28 \\
        14/12 & Confirmação de Performance & \reais{420,00} & 2.00 \\
        \rowcolor{lightGray}
        24-31/12 & Fase de Manutenção (sem intervenção) & \reais{1.140,00} & 1.77 \\
        02/01 & Teste de 8 Variações do Criativo 31 & \reais{280,00} & 1.52 \\
        \rowcolor{lightGray}
        03/01 & Escala CBO + ABO Top 4 & \reais{650,00} & 1.89 \\
        04/01 & Duplicação CBO 131 (orçamento \reais{200,00}) & \reais{866,00} & 2.03 \\
        \bottomrule
    \end{tabularx}
\end{table}

\subsection{Análise por Período}

\textbf{Semana 1 (10-17/12):} Lucro de \reais{1.377,00}
\begin{itemize}
    \item Teste criativo identificou Criativo 31 como vencedor (ROAS 1.8+)
    \item Primeira escala em ABO apresentou oscilação (ROAS 1.28 no dia 12-13)
    \item Paciência validada: ROAS subiu para 2.0 no dia 14
\end{itemize}

\textbf{Semana 2 (17-24/12):} Lucro de \reais{1.000,00}
\begin{itemize}
    \item Fase de consolidação, sem novos testes criativos
    \item ROAS médio: 1.65
\end{itemize}

\textbf{Semana 3 (24-31/12):} Lucro de \reais{2.018,00}
\begin{itemize}
    \item \textbf{Fase de Manutenção:} 7 dias sem intervenção
    \item Campanha ABO 111 (Criativo 31, orçamento \reais{170,00}) entregou ROAS 1.77
    \item ROAS no dia 31/12 (virada de ano): 2.0
\end{itemize}

\textbf{Primeiros 4 Dias/2026 (01-04/01):} Lucro de \reais{2.469,00}
\begin{itemize}
    \item Novo teste criativo (8 variações) validou Criativos 39, 42, 44
    \item Escala em CBO (Criativo 31, \reais{200,00}) performou ROAS 2.0+
    \item Escala em ABO Top 4 (Criativos 14, 31, 42, 44) performou ROAS 1.8+
\end{itemize}

\newpage

\subsection{Lições do Caso Real}

\textbf{1. Oscilação é Normal nos Primeiros Dias de Escala}

No dia 12/12, a campanha de escala entregou ROAS 1.28 (abaixo do ideal). A tentação seria desativar. Mantida ativa, no dia 14/12 o ROAS saltou para 2.0, gerando \reais{420,00} de lucro em um único dia.

\textbf{2. Fase de Manutenção Gera Lucro Passivo}

Do dia 24/12 ao 31/12, nenhuma intervenção foi feita. As campanhas validadas rodaram autonomamente, gerando \reais{2.018,00} de lucro. Isso só é possível quando você:
\begin{itemize}
    \item Valida criativos rigorosamente na Fase 1
    \item Respeita o período de estabilização na Fase 2
    \item Não desativa campanhas prematuramente por oscilações de 1-2 dias
\end{itemize}

\textbf{3. Escala Progressiva vs. Escala Agressiva}

A operação levou 25 dias para passar de \reais{250,00}/dia (teste criativo) para \reais{866,00}/dia (escala consolidada). Tentativas de pular etapas invariavelmente resultam em prejuízo.

\textbf{4. Cardápio de Estruturas}

Ambas ABO e CBO funcionaram. No início (dia 12), ABO performou melhor. Após constância validada (dia 3/01), CBO entregou ROAS superior. Teste ambas.

\subsection{Gráfico de Evolução de Lucro}

\begin{center}
\begin{tikzpicture}
    \begin{axis}[
        width=14cm,
        height=8cm,
        xlabel={Período (Semanas)},
        ylabel={Lucro Acumulado (R\$)},
        xmin=0, xmax=4,
        ymin=0, ymax=7000,
        xtick={1,2,3,4},
        xticklabels={Semana 1,Semana 2,Semana 3,Jan 01-04},
        ytick={0,1000,2000,3000,4000,5000,6000,7000},
        legend pos=north west,
        ymajorgrids=true,
        grid style=dashed,
    ]

    \addplot[
        color=primaryColor,
        mark=*,
        ultra thick
        ]
        coordinates {
        (1,1377)(2,2377)(3,4395)(4,6864)
        };
        \addlegendentry{Lucro Acumulado}

    \end{axis}
\end{tikzpicture}
\end{center}

\textbf{Total em 25 dias:} \reais{6.864,00} de lucro líquido (após desconto de custos tributários e operacionais).

\newpage

\section{Checklist Executivo de Escala}

\subsection{Fase 1: Teste de Criativo (Dias 1-3)}

\begin{itemize}
    \item[$\square$] Estrutura ABO 1-3-1 configurada (1 campanha, 3 conjuntos, 1 anúncio)
    \item[$\square$] Orçamento: 45\%-50\% do ticket médio por conjunto
    \item[$\square$] Público: Advantage+ com sugestões de gênero/idade
    \item[$\square$] Rodou por no mínimo 2 dias completos
    \item[$\square$] ROAS $\geq$ 1.8 confirmado em pelo menos 1 criativo
    \item[$\square$] Métricas secundárias dentro do range (CPV 5\%-6\%, Checkout 25\%-30\%)
\end{itemize}

\subsection{Fase 2: Escala (Dias 4-10)}

\begin{itemize}
    \item[$\square$] Pausou todos os conjuntos com ROAS $<$ 1.3
    \item[$\square$] Manteve em observação conjuntos com ROAS 1.5-1.8
    \item[$\square$] Escolheu estrutura de escala (ABO Modelo 01, ABO Modelo 02 ou CBO)
    \item[$\square$] Orçamento de escala: 3x a 6x o ticket médio
    \item[$\square$] Manteve campanha de escala ativa por no mínimo 3 dias antes de avaliar
    \item[$\square$] ROAS de escala acima de 1.5 (queda aceitável vs. teste criativo)
\end{itemize}

\subsection{Fase 3: Manutenção (Dias 11-25)}

\begin{itemize}
    \item[$\square$] Campanhas performando com ROAS $\geq$ 1.6 por 5+ dias consecutivos
    \item[$\square$] Reduziu intervenções (verificação 1x ao dia ao invés de múltiplas)
    \item[$\square$] Teste semanal de novos criativos (modelagem dos vencedores)
    \item[$\square$] Análise de padrões: identificou estruturas de melhor performance
    \item[$\square$] Duplicação de campanhas vencedoras com orçamento incrementado
\end{itemize}

\subsection{Fase 4: Saturação (Quando Aplicável)}

\begin{itemize}
    \item[$\square$] ROAS em queda constante por 5+ dias consecutivos
    \item[$\square$] CPM em elevação (sinal de saturação de público)
    \item[$\square$] Novos criativos não validam mais (ROAS $<$ 1.5)
    \item[$\square$] Redução progressiva de orçamento antes do encerramento
    \item[$\square$] Preparação de nova oferta seguindo metodologia MVT
\end{itemize}

\newpage

\section{Ferramentas de Análise e Tomada de Decisão}

\subsection{Planilha de Métricas (Template)}

Para cada conjunto de anúncios testado, registre os seguintes KPIs em planilha:

\begin{table}[H]
    \centering
    \footnotesize
    \renewcommand{\arraystretch}{1.2}
    \begin{tabularx}{\textwidth}{X c c c c c c}
        \toprule
        \rowcolor{tableHeader} \textbf{Conjunto} & \textbf{Gasto} & \textbf{Receita} & \textbf{ROAS} & \textbf{CPV} & \textbf{Checkout} & \textbf{CPM} \\
        \midrule
        Criativo 31 - Conj. 1 & 35,00 & 77,00 & 2.20 & 5,2\% & 28\% & 18,50 \\
        \rowcolor{lightGray}
        Criativo 31 - Conj. 2 & 35,00 & 63,00 & 1.80 & 5,8\% & 26\% & 21,00 \\
        Criativo 31 - Conj. 3 & 35,00 & 70,00 & 2.00 & 5,5\% & 27\% & 19,20 \\
        \midrule
        \rowcolor{lightGray}
        \textbf{MÉDIA CRIATIVO 31} & \textbf{105,00} & \textbf{210,00} & \textbf{2.00} & \textbf{5,5\%} & \textbf{27\%} & \textbf{19,57} \\
        \bottomrule
    \end{tabularx}
\end{table}

\textbf{Decisão:} Criativo 31 está VALIDADO (ROAS 2.0, todas as métricas secundárias dentro do range). Levar para escala.

\subsection{Funil de Métricas Secundárias}

\begin{center}
\begin{tikzpicture}[scale=0.9]
    % Impressões
    \draw[fill=primaryColor!20] (0,0) -- (8,0) -- (7,1) -- (1,1) -- cycle;
    \node at (4,0.5) {\textbf{100.000 Impressões (CPM: \reais{18,50})}};

    % Cliques
    \draw[fill=primaryColor!40] (1,1) -- (7,1) -- (6.5,2) -- (1.5,2) -- cycle;
    \node at (4,1.5) {\textbf{3.500 Cliques (CTR: 3,5\%)}};

    % Visualizações de Página
    \draw[fill=primaryColor!60] (1.5,2) -- (6.5,2) -- (6,3) -- (2,3) -- cycle;
    \node at (4,2.5) {\textbf{210 Visualizações (CPV: 5,5\%)}};

    % Checkouts Iniciados
    \draw[fill=primaryColor!80] (2,3) -- (6,3) -- (5.5,4) -- (2.5,4) -- cycle;
    \node at (4,3.5) {\textbf{57 Checkouts (27\%)}};

    % Vendas Aprovadas
    \draw[fill=successGreen!80] (2.5,4) -- (5.5,4) -- (5,5) -- (3,5) -- cycle;
    \node at (4,4.5) {\textbf{20 Vendas (ROAS: 2.0)}};
\end{tikzpicture}
\end{center}

\subsection{Gatilhos de Desativação}

Pause imediatamente um conjunto/campanha quando:

\begin{table}[H]
    \centering
    \renewcommand{\arraystretch}{1.4}
    \begin{tabularx}{\textwidth}{X c}
        \toprule
        \rowcolor{tableHeader} \textbf{Condição} & \textbf{Limite} \\
        \midrule
        ROAS em teste criativo após 3 dias & $<$ 1.3 \\
        \rowcolor{lightGray}
        ROAS em escala após 3 dias & $<$ 1.2 \\
        CPV (Custo por Visualização) & $>$ 10\% \\
        \rowcolor{lightGray}
        CPM (Custo por Mil Impressões) & $>$ \reais{35,00} \\
        Taxa de Checkout & $<$ 15\% \\
        \bottomrule
    \end{tabularx}
\end{table}

\textcolor{accentColor}{\textbf{Exceção:}} Se um criativo teve ROAS excepcional (2.5+) por vários dias e de repente cai para 0.8x em um único dia, aguarde mais 24h antes de pausar (pode ser anomalia temporária do algoritmo).

\newpage

\section{Erros Fatais a Evitar}

\subsection{1. Escalar Baseado em Número de Vendas}

\textbf{Erro:} ``Este criativo vendeu 10 vezes hoje com orçamento de \reais{50,00}. Vou aumentar para \reais{200,00}.''

\textbf{Por que é fatal:} Se o ROAS for 1.2 (prejuízo de 20\%), você está escalando prejuízo. Com orçamento de \reais{200,00}, o prejuízo será 4x maior.

\textbf{Correção:} Escale apenas criativos com ROAS $\geq$ 1.8, independentemente do número de vendas.

\subsection{2. Testar 10+ Criativos Simultaneamente}

\textbf{Erro:} Subir 15 criativos de uma vez para ``acelerar a validação''.

\textbf{Por que é fatal:} Com o imposto de 12,5\%, você pagará taxa sobre \reais{900,00} (15 criativos $\times$ \reais{60,00}) mesmo se nenhum criativo validar. Isso drena capital e cria viés negativo.

\textbf{Correção:} Teste no máximo 3-4 criativos por rodada. Valide, escale, depois teste novos.

\subsection{3. Desativar Campanhas no Primeiro Dia Ruim}

\textbf{Erro:} Campanha de escala deu ROAS 0.9x no primeiro dia. Pausar imediatamente.

\textbf{Por que é fatal:} O algoritmo precisa de 2-3 dias para otimizar. Você pode estar interrompendo uma campanha que seria lucrativa.

\textbf{Correção:} Aguarde 3 dias completos. Só pause se ROAS acumulado for $<$ 0.7x.

\subsection{4. Copiar Ofertas Saturadas}

\textbf{Erro:} Ver oferta bem-sucedida na Biblioteca de Anúncios, clonar, e esperar o mesmo resultado.

\textbf{Por que é fatal:} Você está competindo com o original + dezenas de outros clonadores. CPM inflado + tributação = prejuízo garantido.

\textbf{Correção:} Use metodologia MVT. Modele princípios, não ofertas. Crie concepções únicas.

\subsection{5. Ignorar Métricas Secundárias}

\textbf{Erro:} ``ROAS está 2.0, não preciso olhar CPV, Checkout, CPM.''

\textbf{Por que é fatal:} Criativos com ROAS alto mas métricas secundárias ruins saturam rapidamente. Você escala, e em 5 dias o ROAS desaba.

\textbf{Correção:} Valide ROAS \textbf{e} métricas secundárias. Priorize criativos que atendam todos os critérios.

\newpage

\section{Considerações Finais}

O framework apresentado neste documento não é teórico. É o resultado de operações reais executadas em dezembro/2025 e janeiro/2026, sob as novas condições tributárias do Facebook Ads.

\subsection{Três Princípios Inegociáveis}

\begin{enumerate}
    \item \textbf{Validação antes de Escala:} Jamais aumente orçamento em campanhas com ROAS $<$ 1.8. A tentação de ``forçar'' uma oferta ruim a funcionar resulta em prejuízo proporcional ao investimento.

    \item \textbf{Paciência Estratégica:} Campanhas de escala oscilam nos primeiros 3 dias. Operadores impacientes desativam cedo e nunca atingem fase de manutenção (lucro passivo).

    \item \textbf{Vantagem Competitiva:} Com a tributação de 12,5\%, apenas ofertas com posicionamento único (metodologia MVT) sustentam margens saudáveis a longo prazo.
\end{enumerate}

\subsection{Próximos Passos}

Após dominar este framework:

\begin{itemize}
    \item \textbf{Semana 1-2:} Execute teste criativo com 3 criativos (estrutura ABO 1-3-1)
    \item \textbf{Semana 3:} Escale o(s) criativo(s) validado(s) usando ABO ou CBO
    \item \textbf{Semana 4:} Entre em fase de manutenção, teste novos criativos semanalmente
    \item \textbf{Mês 2:} Chegue a investimento diário de \reais{300,00}-\reais{500,00} com ROAS $\geq$ 1.6
\end{itemize}

\subsection{Recursos Complementares}

Para aprofundar a metodologia MVT (Modelagem, Vantagem, Tráfego), incluindo:
\begin{itemize}
    \item Frameworks de concepção de oferta
    \item Templates de página de vendas
    \item Análise de criativos vencedores
    \item Acompanhamento ao vivo 2x por semana
\end{itemize}

Entre em contato via WhatsApp ou Instagram para conhecer a \textbf{Mentoria Ou Processo}.

\vspace{2cm}

\begin{center}
\colorbox{lightGray}{%
    \parbox{0.85\textwidth}{%
        \centering
        \vspace{0.3cm}
        \textcolor{primaryColor}{\textbf{``Não se copia ofertas. Cria-se concepções únicas,}}\\
        \textcolor{primaryColor}{\textbf{surfa-se a onda antes da concorrência,}}\\
        \textcolor{primaryColor}{\textbf{e coloca-se lucro no bolso.''}}\\
        \vspace{0.3cm}
    }%
}
\end{center}

\end{document}
